\chapter{Protocol for the two-person study}
\label{app:study}
\label{ch:study}

\noindent\emph{This is the study \cref{ch:methodsmain} says the system was
built for. The pilot in \cref{app:pilot} is a smaller, operator-only
precursor; this protocol -- with a separate wearer -- is still the target.}

\section{The question and hypotheses}

Does supplying part of the operator's control with autonomy improve
performance, and what does it cost the person carrying the arms? No prior
study measures both under the same autonomy in the same trial
(\cref{ch:background}). Five hypotheses, stated with direction before any
data exists: task completion time is shorter under shared autonomy (H1);
operator workload is lower (H2); \textbf{wearer-reported trust is lower
under shared autonomy} (H3), against the direction of H1/H2, taken from a
concealed-operator study~\cite{zhou2026proxemics} and tested here with
autonomy that is genuinely a program and a wearer with a concurrent task;
the wearer's sense of agency is lower while ownership is unchanged (H4),
since the two dissociate~\cite{longo2008embodiment}; cross-arm interference
is lower under shared autonomy while phase lag is not (H5), which would
identify the benefit as attentional. \textbf{H3 and H4 are the hypotheses
worth running the study for}; H1 and H2 exist to establish the platform can
show an effect at all.

\section{Design}

Two modes (direct teleoperation, shared autonomy) crossed with five tasks,
within subjects, sixteen dyads, roles fixed for the whole session. Two modes
rather than four is costed, not a convenience: four modes across five tasks
needs \SIrange{260}{320}{\minute} against a wearer exposure cap of
\SI{48}{\minute} for carrying \SI{17}{\kilo\gram} continuously; a
between-subjects design across modes needs ninety-six participants for the
same sensitivity, since between-subject variance here exceeds the effect
sought. Mode order is counterbalanced with Williams squares, not a Latin
square, because a Latin square balances position but not immediate sequence,
and a task with a learning curve is sensitive to the sequence.

\textbf{Grasping under direct teleoperation is not achievable}: the default
orientation mode pins the wrist to the anchor, the pick task needs a
\SI{169.7}{\degree} rotation from it, and no master channel can command
that. T2 (pick and place) therefore carries no mode contrast and is analysed
as a single condition -- retained because grasp success under autonomy is
worth measuring, and because it is the one task where autonomy's advantage
over teleoperation on this platform is categorical rather than quantitative.

\begin{table}[htbp]
  \centering
  \small
  \caption[Task set]{}
  \label{tab:taskset}
  \begin{tabular}{@{}p{0.20\linewidth}p{0.48\linewidth}p{0.24\linewidth}@{}}
    \toprule
    Task & Measures & Mode contrast \\
    \midrule
    T1 positioning & Throughput vs.\ index of difficulty~\cite{fitts1954information}; single- vs.\ both-arm decrement & both modes \\
    T2 pick and place & Grasp success, placement accuracy & shared autonomy only \\
    T3 rigid carry & Coordination error as object tilt & both modes \\
    T4 compliant carry & Coordination error as gripper separation & both modes \\
    T5 dual pursuit & Cross-arm interference under divided attention & both modes \\
    \bottomrule
  \end{tabular}
\end{table}

\section{Measures and procedure}

NASA-TLX~\cite{hart1988nasatlx} after every block, both participants
separately (the two roles load differently, and the comparison between them
is the object of the study); Borg CR10~\cite{borg1982perceived} exertion from
the wearer between every block, since a single end-of-session rating turns a
covariate into a guess; trust and embodiment~\cite{longo2008embodiment} at
the end of session, wearer only for embodiment. A safety observer, present
throughout, holds an emergency stop and watches the wearer rather than the
screen; their presence is enforced in software as well as procedurally --
the safety node treats two seconds of heartbeat silence as withdrawal of
consent to move. Consent is taken separately, in separate rooms.

\section{Ethics and risk}

The dominant risk is a \SI{17}{\kilo\gram} manipulator pair, moving under
partly autonomous control, within reach of a seated person's head and
torso. Mitigations: the \SI{150}{\milli\metre} margin enforced over the
whole path and measured geometrically (\cref{sec:clearance}); a stop within
the wearer's own reach and a second held by the observer; twenty minutes of
carried-load exposure at a time with exertion rated between blocks; order
counterbalanced against fatigue confounding condition; consent restated at
every break; identifying data refused in code
(\repo{write\_manifest()} raises on name, email, date of birth, address or
telephone number). Video recording is not part of the protocol: a video of a
participant wearing this rig is not anonymous, and no measure here requires
one.
